% !TeX root = ../../Book1.tex
% 纲要标题：### 19.3 显式计算
% 纲要位置：计划纲要.md，第 660 行。

\section{显式计算}
\label{b1:sec:1903}

% 纲要内容（仅作写作计划，尚非教材正文）：
%
% * 二次、多二次及三次扩张
% * 四次方程的典型 Galois 群
% * 从置换作用和分解型识别子群
% * 模素数分解与循环型：在此完整证明约化判据，第21.3节直接引用
%

识别 Galois 群通常需要两类信息：扩张次数确定群阶，群对根的作用将它表示为对称群的子群。以下每个例子都先明确分裂域，再计算自同构；不能只写出一个抽象群名而略去它怎样作用。

\subsection{二次、多二次与三次扩张}
\label{b1:subsec:1903:01}

特征不为二时，若 \(d\in K^\times\) 不是平方，\(K(\sqrt d)/K\) 为二次 Galois 扩张，非平凡自同构把 \(\sqrt d\) 送到 \(-\sqrt d\)。更一般的可分二次多项式 \(x^2+bx+c\) 的两根之和为 \(-b\)，所以添加一根后另一根已在域内；即使特征为二，可分二次扩张也仍为 Galois 扩张。

\ProseParagraphBreak
考虑 \(L=\mathbb Q(\sqrt2,\sqrt3)\)。若 \(\sqrt3=a+b\sqrt2\)，其中 \(a,b\in\mathbb Q\)，平方比较得 \(2ab=0\)、\(a^2+2b^2=3\)。若 \(b=0\)，三为有理数平方；若 \(a=0\)，\(3/2\) 为有理数平方；两者都与素因数指数的奇偶性矛盾。因此 \([L:\mathbb Q]=4\)，基可取 \(1,\sqrt2,\sqrt3,\sqrt6\)。独立改变两个平方根的符号保持全部乘法关系，给出四个自同构，所以
\[
\operatorname{Gal}(L/\mathbb Q)\cong C_2\times C_2.
\]
三个二元子群的不动域分别为 \(\mathbb Q(\sqrt2)\)、\(\mathbb Q(\sqrt3)\)、\(\mathbb Q(\sqrt6)\)。例如同时改变两者符号的自同构固定 \(\sqrt6\)，而不动域次数为二，故恰为最后一个域。

\ProseParagraphBreak
对 \(x^3-2\)，令 \(a=\sqrt[3]2\in\mathbb R\)、\(\zeta^2+\zeta+1=0\)，取非实的 \(\zeta\)。分裂域为 \(L=\mathbb Q(a,\zeta)\)。Eisenstein 判据给出 \([\mathbb Q(a):\mathbb Q]=3\)，而 \(\zeta\) 不在实域 \(\mathbb Q(a)\) 中，所以总次数为六。自同构忠实地置换三个根 \(a,\zeta a,\zeta^2a\)，因此是 \(S_3\) 的一个六元子群，只能为 \(S_3\)。其中复共轭固定 \(a\)，并交换另外两个根。

\subsection{四次方程的 Galois 群}
\label{b1:subsec:1903:02}

四次不可约多项式的分裂域次数未必为四。三个典型计算分别给出不同的群。

\ProseParagraphBreak
先取 \(f=x^4-10x^2+1\)。它的根为 \(\pm\sqrt2\pm\sqrt3\)。令 \(a=\sqrt2+\sqrt3\)，则 \(a^{-1}=\sqrt3-\sqrt2\)，所以
\[
\sqrt3=\frac{a+a^{-1}}2,\qquad \sqrt2=\frac{a-a^{-1}}2.
\]
由上一小节的次数计算，\([\mathbb Q(a):\mathbb Q]=4\)。因此 \(f\) 是 \(a\) 的最小多项式，且其分裂域就是 \(\mathbb Q(\sqrt2,\sqrt3)\)，Galois 群为四元群 \(C_2\times C_2\)。

\ProseParagraphBreak
再取 \(f=x^4+ x^3+x^2+x+1\)，其根为非一的五次单位根。变元平移后的 \(f(x+1)\) 由素数五满足\cref{b1:thm:eisenstein}，故 \(f\) 不可约。若 \(\zeta\) 是一个根，所有根 \(\zeta,\zeta^2,\zeta^3,\zeta^4\) 都在 \(\mathbb Q(\zeta)\) 中，扩张次数为四。自同构由 \(\zeta\mapsto\zeta^j\) 给出，复合对应指数模五相乘。模五的二依次给出 \(2,4,3,1\)，所以这个 Galois 群为 \(C_4\)。

\ProseParagraphBreak
最后取 \(f=x^4-2\)，令 \(a=\sqrt[4]2>0\)。分裂域为 \(L=\mathbb Q(a,i)\)。Eisenstein 判据给出实子域 \(\mathbb Q(a)\) 的次数为四；加入 \(i\) 后次数变为八。设 \(r\) 固定 \(i\) 而把 \(a\) 送到 \(ia\)，设 \(s\) 为复共轭。前者确实存在，因为 \([L:\mathbb Q(i)]=4\)，\(a\) 在 \(\mathbb Q(i)\) 上的最小多项式仍是 \(x^4-2\)，可以把 \(a\) 映到另一个根 \(ia\)。直接作用于生成元可验
\[
r^4=s^2=1,\qquad srs=r^{-1}.
\]
四个 \(r^j\) 固定 \(i\)，四个 \(sr^j\) 把 \(i\) 变号，彼此不同，已经列出全部八个自同构。因此群为本书记号下的 \(D_8\)。

\subsection{置换作用、分解型与子群识别}
\label{b1:subsec:1903:03}

若首一可分多项式 \(f\in K[x]\) 的分裂域为 \(L\)，其 Galois 群在各根上的作用忠实，因为固定全部根就固定它们生成的域。一个不可约因子的根构成一个轨道：任意两根由简单扩张同构对应，该同构能由\cref{b1:thm:embedding-extension}延拓到 \(L\)，正规性保证延拓为自同构。因此在 \(K\) 上各不可约因子的次数，正是全部根分成的轨道大小。特别地，不可约多项式给出传递作用。

\ProseParagraphBreak
还可利用判别式辨认置换的奇偶性。设特征不为二，互异的根为 \(a_1,\ldots,a_n\)，并置 \(\delta=\prod_{i<j}(a_j-a_i)\)。交换两个根使 \(\delta\) 变号，因此任意 \(\sigma\in G\) 满足
\[
\sigma(\delta)=\operatorname{sgn}(\sigma)\delta.
\]
\(\delta^2\) 是多项式判别式，且被 \(G\) 固定，故属于 \(K\)。它是 \(K\) 中的平方当且仅当 \(\delta\in K\)：若 \(\delta^2=c^2\)，域中有 \(\delta=\pm c\)；反向显然。再由 \(L^G=K\)，得到
\begin{equation}\label{b1:eq:discriminant-alternating-group}
G\subseteq A_n\quad\Longleftrightarrow\quad
\operatorname{disc}(f)\text{ 是 }K\text{ 中的平方}.
\end{equation}
例如一个可分不可约三次多项式的 Galois 群是 \(S_3\) 的传递子群，其阶被三整除且整除六，所以只有 \(A_3\cong C_3\) 与 \(S_3\) 两种可能；判别式是否为平方将它们区分。

\ProseParagraphBreak
这里的轨道大小取自原基域上的不可约分解。若改在有限域中分解整数多项式，则可以借助 Frobenius 自同构得到原 Galois 群中单个元素的循环型。下面建立两者的联系。

\subsection{模素数分解与循环型}
\label{b1:subsec:1903:04}

对一个整数多项式，模素数后的因子次数往往容易计算。要把这些数据用于特征零中的 Galois 群，须使有限域中置换根的 Frobenius 自同构来自原分裂域的自同构。以下证明用根生成的整数环及其有限商域实现这一点；用到的工具是\cref{b1:lem:integral-closure-elementary}的整性封闭性、\cref{b1:thm:crt}的中国剩余定理，以及已经建立的有限域和有限 Galois 理论。初读只关注 Galois 对应的读者，可以在需要计算具体多项式的群时再读本小节。

\begin{lemma}\label{b1:lem:reduction-cycle-type}
设 \(f\in\mathbb Z[X]\) 首一，在 \(\mathbb Q\) 上无重根；素数 \(p\) 满足模 \(p\) 的约化 \(\bar f\) 也无重根。若 \(\bar f\) 的不同首一不可约因子的次数为 \(d_1,\ldots,d_r\)，则 \(f\) 在 \(\mathbb Q\) 上的 Galois 群，在全部根上的置换作用中包含循环长度为 \(d_1,\ldots,d_r\) 的元素。
\end{lemma}
\begin{proof}
设分裂域为 \(L\)，根为 \(\alpha_1,\ldots,\alpha_n\)，\(G=\operatorname{Gal}(L/\mathbb Q)\)，并取环 \(A=\mathbb Z[\alpha_1,\ldots,\alpha_n]\)。所有根都满足首一整数方程，\cref{b1:lem:integral-closure-elementary}说明 \(A\) 由有限多个元素作为 \(\mathbb Z\) 模生成，且每个元素都在 \(\mathbb Z\) 上整。\(G\) 置换这些根，所以保持 \(A\)。由 \(L^G=\mathbb Q\)，\(A^G\) 中的元素为有理整元素，只能是整数：一个既约分数若满足首一整数方程，清除分母后就要求分母整除分子的某次幂，故分母为一。因此 \(A^G=\mathbb Z\)。

理想 \(pA\) 为真理想，否则 \(1/p\in A\)，与刚才的有理整元素判别矛盾。由\cref{b1:prop:maximal-ideal-exists}取极大理想 \(\mathfrak m\supseteq pA\)。商域 \(k=A/\mathfrak m\) 有限，因为 \(A/pA\) 由有限多个元素作为 \(\mathbb F_p\) 向量空间生成，而 \(k\) 是它的商。令
\[
D=\{\,\sigma\in G\mid\sigma(\mathfrak m)=\mathfrak m\,\},
\]
这是 \(\mathfrak m\) 在 \(G\) 作用下的稳定子群，故为子群；它在 \(k\) 上诱导自同构，记其像为 \(H\leq\operatorname{Gal}(k/\mathbb F_p)\)。我们证明这个像为全群。

任取 \(x\in k^H\)。\(G\) 作用下 \(\mathfrak m\) 的轨道由有限个不同的极大理想组成，它们两两互素。\cref{b1:thm:crt}给出 \(a\in A\)，在 \(\mathfrak m\) 处的余数为 \(x\)，在其余轨道理想处的余数为零。多项式
\[
P(T)=\prod_{\sigma\in G}\Bigl(T-\sigma(a)\Bigr)
\]
的系数属于 \(A^G=\mathbb Z\)。若 \(\sigma\in D\)，因 \(x\) 被 \(H\) 固定，\(\sigma(a)\) 模 \(\mathfrak m\) 为 \(x\)；若 \(\sigma\notin D\)，则 \(a\) 模 \(\sigma^{-1}(\mathfrak m)\) 为零，故 \(\sigma(a)\) 模 \(\mathfrak m\) 为零。写 \(|D|=p^e u\)、\(p\nmid u\)，约化后得到
\[
\bar P(T)=T^{|G|-|D|}(T-x)^{|D|}
=T^{|G|-|D|}\Bigl(T^{p^e}-x^{p^e}\Bigr)^u.
\]
其中 \(T^{|G|-p^e}\) 的系数为 \(-u x^{p^e}\)，但 \(\bar P\) 的全部系数在 \(\mathbb F_p\) 中。由于 \(u\) 在该域中可逆，\(x^{p^e}=c\in\mathbb F_p\)，于是 \((x-c)^{p^e}=0\)，得到 \(x=c\)。所以 \(k^H=\mathbb F_p\)。有限域扩张 \(k/\mathbb F_p\) 为 Galois 扩张，\cref{b1:thm:finite-galois-correspondence}迫使 \(H=\operatorname{Gal}(k/\mathbb F_p)\)。

因此 Frobenius \(b\mapsto b^p\) 在 \(D\) 中有一个提升 \(\sigma\)。等式 \(f(X)=\prod_i(X-\alpha_i)\) 约化后表明 \(\alpha_i\) 的像是 \(\bar f\) 的全部根。由于约化无重根，这些像互不相同。由\cref{b1:thm:finite-field-subfields}及\cref{b1:thm:finite-field-automorphisms}，一个次数为 \(d_j\) 的不可约因子的根恰组成长度 \(d_j\) 的轨道：根生成的域有 \(p^{d_j}\) 个元素，而固定该根的最小正 Frobenius 次幂正是 \(d_j\)。提升 \(\sigma\) 在原根集上的循环类型与约化后的 Frobenius 相同，遂得结论。
\end{proof}

因此，每次无重根的模素数分解都给出群中的一种循环型。将不同素数提供的信息与传递性、判别式判据结合，往往能确定整个群。\cref{b1:prop:quintic-x5-x-1}将用这一方法计算 \(X^5-X-1\) 的 Galois 群；这里的工具本身不依赖根式可解性理论。
